\section{Experimental Integrity and Deviations}
\label{sec:integrity}

\paragraph{What was frozen, and when.} The experiment specification --- factors
and levels, criteria, seed count, the resolution criterion, screening gates,
baseline ladder, model class and hyperparameters, abstention rules and shift
splits --- was written and committed to version control \emph{before any
evaluation run was executed}. Policy parameters were frozen at the same time and
varied afterwards only as sensitivity (Section~\ref{sec:sensitivity}). No factor level was chosen
because it produced a particular outcome; the demand levels follow from measured
capacity.

\paragraph{Execution reconciliation.} The main campaign comprises
\NumRunsMain{} valid runs ($\NumContexts \times \NumPolicies \times \NumSeeds$),
and this is the number underlying every result in Section~5 unless stated
otherwise. The project executed \NumRunsTotalAttempted{} simulations in total,
of which \NumRunsTotalSucceeded{} succeeded: \NumRunsValidate{}
scenario-validation, \NumRunsScreen{} screening, \NumRunsMain{} main,
\NumRunsSens{} sensitivity, \NumRunsOODNorth{} domain-shift and
\NumRunsBoundary{} boundary-refinement runs, plus \NumRunsFailedBypass{} runs
that failed for the reason given below. These counts are distinct and are not
interchangeable.

\paragraph{Protocol deviation D-1.} The pre-specified complementarity screen did
not satisfy criterion G1 and therefore triggered the declared stopping rule. The
main campaign was nevertheless executed, under the documented protocol deviation
D-1. The reasons for continuation were recorded before campaign launch: a
selection problem demonstrably existed; G1 was written against the scalar
criterion and was blind to the three-criterion structure beside it; and the
negative finding rested on 16 of \NumContexts{} contexts sampled at extreme
factor levels only. A confirmatory criterion G1$'$, declared before the main
campaign and identical in definition to G1, was then evaluated on all
\NumContexts{} contexts. G1$'$ \NumGOneVerdict{}
(Section~\ref{sec:complementarity}).

We do not present G1$'$ as erasing the deviation, and we do not describe the
original screen as having passed. It did not. The methodological lesson is that
the screen sampled corner conditions and did not sample the interior transition
region, so it had power to detect near-ties but not coverage to find where the
minority policies win. A resolution-IV fraction at extreme levels is an
efficient design for main effects and a poor one for discovering where a winner
map changes. That is a limitation of our screening design, not a justification
for having continued.

\paragraph{An infeasible shift split and its replacement.} One declared shift
split placed the disruption on the bypass corridor. The bypass is single-lane,
so the specified capacity-reducing lane closure does not constrain it --- it
disconnects it, leaving vehicles already routed there without a valid path. All
\NumRunsFailedBypass{} attempted runs failed for this reason. The split was
replaced by the nearest feasible construction, a disruption on the two-lane
north corridor, which preserves the intent (a corridor carrying no disruption in
any training context) while remaining a capacity reduction rather than a
severance. A guard now refuses any incident on a single-lane edge. The
replacement is reported as split O8 and is \emph{not} the originally declared
split.

\paragraph{Traceability.} Every numeric claim in this paper is generated from a
single machine-readable results file by a script, and rendered through a macro;
a build-time check fails if the text references a quantity the analysis did not
produce. Seventeen automated quality-control checks cover leakage discipline,
training-fold-only baseline selection, freeze ordering, shift reporting and
citation completeness. All \NumRefs{} references were verified against publisher
records; DOIs that could not be confirmed directly were omitted rather than
inferred.
